Generative AI was used in two bounded and disclosed parts of the project. First, the repository labels the GitHub Actions workflow for TypeDoc deployment as AI-generated. Second, AI assistance was used to inspect the repository, verify certain sections of the report, revise the introduction and background for consistency, add source references, and check whether the report's claims matched the implementation.

AI output was treated as a draft rather than evidence. Security statements were checked against the TypeScript, Kotlin, SQL, Docker, and Gradle files, and external recommendations were tied to the cited OWASP, Node.js, CISA, and Android documentation. This verification changed the report materially: it prevented claims that Slowloris protection and dynamic attack measurements already existed, and it exposed the mismatch between opaque login tokens and JWT validation, the cleartext development transport, and the missing \verb|await| during registration.

No secret values from the ignored local environment file were copied into the report. The team remains responsible for reviewing the generated prose, confirming experimental results when a runnable environment is available, and ensuring that the final submission complies with the course's policy on AI-assisted work.
